\chapter{Streaming Pipeline and Integration}
\label{ch:streaming}

This chapter describes the runtime pipeline: how events flow from ingestion through enrichment, scoring, and decision in near real time, the software pattern that enables code reuse across streaming, batch, and API paths, the fault tolerance layer, and the integration with PostgreSQL and the nightly rebuild.

\section{Topology and Broker}

The pipeline is a linear chain of four topics, each consumed by a dedicated microservice producing to the next, all sharing \code{client\_id} as partition key for the ordering reason established in Chapter~\ref{ch:architecture}.

RedPanda provides a Kafka-compatible API without the JVM or ZooKeeper dependency, which matters in a Docker Compose deployment where every service competes for memory: a Kafka cluster needs at least three containers, while RedPanda runs as one with the same producer/consumer API. The streaming wrappers use \code{kafka-python} and are unaware they are talking to RedPanda rather than Kafka --- the switch is transparent at the application level.

Each service registers a dedicated consumer group, which provides two guarantees the pipeline relies on. \emph{Offset tracking}: a restarted service resumes from its last committed offset rather than replaying the topic. \emph{Horizontal scalability}: adding a second instance automatically splits partitions across the two consumers, doubling throughput without code changes. The current deployment runs one consumer per service, but the architecture supports scaling out by adding replicas.

\section{Core Extraction Pattern}

The system supports three execution paths that all require the same computation: streaming (Kafka loop, one event at a time, writes to Redis buffers), batch (Airflow DAG over PostgreSQL rows, no Redis), and the simulation API (synchronous request--response, read-only Redis). Duplicating enrichment, scoring, or decision logic across them would create both a maintenance burden and a drift risk.

Each stage therefore has a \textbf{core module} encapsulating all domain logic: \code{EnrichmentCore.enrich\_event()}, \code{ScoringCore.score\_event()}, and \code{DecisionCore.decide\_event()}. Each core receives its dependencies through constructor injection --- a Redis client, model objects, configuration --- exposes a single public method, and has no knowledge of Kafka, Airflow, or FastAPI: it operates on Python dictionaries and returns Python dictionaries.

The Kafka-specific logic lives in thin wrappers that do exactly three things: poll a topic, call the core's public method, and produce the result. The \code{\_\_main\_\_} block is the sole wiring point --- it creates the Redis client, loads models, constructs the core, injects it into the wrapper, and starts the loop.

The payoff was concrete. When the batch pipeline was rewritten it imported \code{EnrichmentCore} and called the same \code{\_compute\_features()} method; when the simulation API was built it imported all three cores and wired them with a shared Redis client. In both cases zero feature computation code was duplicated.

\section{Dead Letter Queue and Fault Tolerance}

In a streaming pipeline a single malformed event can block an entire consumer. The fault tolerance layer classifies failures into two categories with opposite responses.

\textbf{Poison pills} (\code{KeyError}, \code{ValueError}, \code{JSONDecodeError}) are structurally invalid and will never succeed regardless of retries; they are routed to the DLQ immediately with \code{retry\_count=0}. \textbf{Transient failures} (\code{redis.ConnectionError}, \code{psycopg2.OperationalError}, network timeouts) may resolve on their own; the consumer applies exponential backoff (1\,s, 2\,s, 4\,s) and retries up to three times before routing to the DLQ.

The distinction matters because the two classes demand opposite handling: retrying a poison pill wastes time and will never succeed, while immediately discarding a transient failure loses an event that would have succeeded one second later when Redis reconnected.

A \code{ResilientConsumer} class wraps all three streaming services. It intercepts exceptions from the core's public method, classifies them with an \code{isinstance} check against a \code{TRANSIENT\_ERRORS} tuple, applies the appropriate strategy, and produces failed events to a single \code{dlq} topic wrapped in a metadata envelope: the original event, source topic, failing service, error type and message, full stacktrace, retry count and maximum, and first/last failure timestamps. The envelope preserves enough context for an operator to diagnose the failure, replay the event after fixing the root cause, and distinguish infrastructure outages (transient, clustered in time) from data quality issues (poison pills, scattered across clients).

\section{Enrichment Service}

The Enrichment Service consumes raw events, computes the 30 contrast features by comparing the current event against the client's historical profile, and produces enriched events.

For each event it fetches three items from Redis: the client profile (\code{profile:client:\{id\}}, batch-computed statistics --- amount means and stds per operation, operation mix distributions, archetype, account age), the client's rolling buffer of the last 50 raw event summaries for velocity features, and the employee buffer for peer comparison features. The profile is rebuilt nightly by the batch pipeline; the buffers are updated per event by the service itself after feature computation.

The term \emph{contrast feature} is central: raw event fields carry no anomaly signal on their own, since a 10{,}000~DT versement is normal for a big business and extreme for a student. Every feature is computed relative to the client's own baseline, turning absolute values into deviations comparable across archetypes.

Two implementation details deserve mention. When a client has no 30-day history for a given operation type --- common for infrequent operations like cheque deposits --- the z-score computation would divide by zero. A three-level \textbf{fallback chain} applies: 30-day statistics, then 90-day, then the archetype baseline, ensuring every event receives a meaningful contrast score regardless of transaction density. And after computing the features, the service appends the \emph{unscaled} vector to \code{sequence:client:\{id\}}, capped at 10 entries; the scaler is owned by the Scoring Service, and storing pre-scaled features would couple the Enrichment Service to the scaler version, breaking independence during retraining.

\section{Scoring Service}

The Scoring Service runs both models and produces a fused score with optional SHAP explanations. The Autoencoder's reconstruction error and the LSTM's prediction error are each converted to a percentile rank against the training distribution, then combined with a sigmoid weight controlled by account maturity:
\[
w_{\text{lstm}} = \sigma\!\left(\frac{d - 90}{15}\right), \qquad
\text{fused} = (1 - w_{\text{lstm}}) \cdot s_{\text{ae}} + w_{\text{lstm}} \cdot s_{\text{lstm}}
\]
where $d$ is days since account opening. For accounts younger than $\sim$75 days the Autoencoder dominates; beyond $\sim$105 days the LSTM does. The 90-day midpoint reflects the minimum history for the sequence model to produce reliable predictions.

\textbf{Cold-start guard.} Clients with fewer than 5 events in the sequence buffer pose a distinct problem: their velocity features sit far outside the training distribution, which applied a 21-day warm-up filter. Without a guard, these zero-velocity vectors produce reconstruction errors mapping above the 0.99 percentile, generating a false BLOCK for every new client. The service therefore checks buffer length and returns a neutral 0.5 with $w_{\text{lstm}} = 0$ and no SHAP when the buffer holds fewer than 5 events, transitioning to AE-only at 5--9 and full fusion at 10+.

\textbf{Conditional SHAP.} Because KernelExplainer requires multiple model evaluations per event, explanations are generated only when the fused score exceeds $T_1 = 0.95$. Events below that threshold are overwhelmingly normal and their explanations would never reach a supervisor's screen.

\section{Decision Service}

The Decision Service consumes scored events and produces the final tier --- INFO, REVIEW, ALERT, or BLOCK --- with an explanation combining model scores and rule triggers. The base tier comes from the three score thresholds of Chapter~\ref{ch:architecture}; the eight regulatory rules then evaluate, and each can \emph{escalate} the tier but never de-escalate it.

The escalation-only principle exists because rules encode hard regulatory requirements that must never be overridden by a statistical model: a model might assign a low score to a transaction that happens to hit a reporting threshold, and the rule ensures the supervisor still sees it. Regulatory flags are logged for audit whether or not they change the tier.

Events flagged \code{cold\_start=True} still pass through all rules --- a genuinely suspicious transaction from a new client must trigger the appropriate flags --- but the final tier is \textbf{capped at REVIEW}. This prevents the neutral 0.5 score from being escalated to ALERT or BLOCK by rules alone, which would generate high-priority alerts for clients who simply have not built enough history yet. The cap is recorded in the decision reasons so supervisors understand the limit.

\section{Sink Consumers and Persistence}

Two sink consumers persist streaming data to PostgreSQL. The \code{raw\_events\_sink} inserts every event into the \code{transactions} table, providing the permanent record used by the nightly profile rebuild and the dashboard's transaction history. The \code{decisions\_sink} performs two writes per event: an insert into the \code{alerts} table with tier, fused score, SHAP explanation as JSONB, and triggered rule names (with \code{supervisor\_decision} initially NULL); and an increment of the client's risk-history counters in Redis, which feed back into enrichment features for subsequent events.

\subsection{Deliberate FK Omission}

The \code{alerts} table has \textbf{no foreign key} to \code{transactions} despite sharing \code{event\_id}. This is intentional. The two sinks run in parallel, and because the decision path adds latency (enrichment, scoring, rule evaluation), the decisions sink may attempt to insert an alert before the raw events sink has inserted the corresponding transaction. A foreign key would cause intermittent insertion failures depending on which consumer happens to be faster. The \code{event\_id} still links the tables logically; the constraint is enforced at the application level rather than the database level.

\subsection{Schema Overview}

PostgreSQL is the cold store for the whole system, written by the two sinks and the batch pipeline, and read by the dashboard, the profile rebuild, and the retraining pipeline.

\begin{table}[htbp]
\centering
\small
\caption{PostgreSQL tables and their roles}
\begin{tabular}{@{}p{4.1cm}p{2.9cm}p{6.2cm}@{}}
\toprule
\textbf{Table} & \textbf{Key} & \textbf{Role} \\
\midrule
\code{clients\_master} & \code{client\_id} & Archetype, account opening date, client type, home branch \\
\code{employees\_master} & \code{employee\_id} & Teller identity and assigned branch \\
\code{accounts} & \code{account\_id} & Account reference data (FK to client) \\
\code{transactions} & \code{event\_id} & All raw events; written by the raw events sink \\
\code{evaluation\_labels} & \code{event\_id} (FK) & Ground truth: oracle labels and supervisor verdicts \\
\code{alerts} & \code{alert\_id} & Tier (ENUM), fused score, reasons, regulatory flags, JSONB explanation, supervisor decision \\
\code{profile\_snapshots} & (client, date) & Nightly client profile as JSONB \\
\code{employee\_profile\_snapshots} & (employee, date) & Nightly employee profile as JSONB \\
\code{archetype\_baselines} & \code{archetype} & Per-archetype reference statistics (final fallback) \\
\bottomrule
\end{tabular}
\end{table}

Three schema decisions are worth noting. The \code{tier} column uses a PostgreSQL \code{ENUM} rather than \code{VARCHAR}, enforcing valid values at the database level and enabling direct ordering (INFO $<$ REVIEW $<$ ALERT $<$ BLOCK) without a \code{CASE} expression. Oracle labels live in a separate \code{evaluation\_labels} table because they are generator metadata that does not exist in a production system --- separating them makes exclusion trivial without touching the transaction schema. And profiles are stored as \textbf{JSONB} rather than normalized columns: a client profile has $\sim$194 fields across 6 categories, and normalizing them would produce a table whose shape changes whenever a feature is added. The tradeoff is that individual fields cannot be indexed efficiently, which the access pattern (load the whole profile into Redis) does not require.

\section{Batch Integration}

The streaming pipeline depends on client profiles too expensive to recompute per event. A nightly Airflow DAG (\code{nightly\_profiles}) rebuilds them from the full transaction history through eight tasks with validation gates: rebuild client snapshots, validate completeness, rebuild employee snapshots, validate, recompute archetype baselines, validate, hydrate Redis, and run a smoke test scoring a sample event with the new profiles. Every task uses \code{ON CONFLICT} upserts on its target table, so any task can be retried after a mid-execution failure without duplicating or corrupting rows.

\textbf{Feature parity guarantee.} The batch computation calls the same \code{EnrichmentCore.\_compute\_features()} method as the streaming path. Because the method is pure --- taking event, profile, and buffer data as arguments and returning 30 features with no side effects --- the only difference is that the batch script passes the archetype baseline as a parameter instead of fetching it from Redis. Any change to feature computation is therefore automatically reflected in both paths.

\textbf{Streaming--batch coordination} is a known limitation. The current deployment does not pause consumers during the rebuild window; in practice the DAG runs during low-traffic hours and the hydration step overwrites profile keys atomically, leaving a brief window in which consecutive events could read pre- and post-rebuild profiles. This one-event inconsistency is acceptable for the current deployment; a production system would pause consumers during hydration or swap profile versions inside a Redis transaction. This is documented as future work in Chapter~\ref{ch:roadmap}.
